\subsection{Global place balancing: an exact lattice and a stable extension
theorem}

The preceding local calculations determine all coefficients that occur in the
source's scalar component formula, but they do not yet make that formula a
connecting map of one global extension.  This distinction produces a concrete
construction problem rather than a slogan.

Fix an ordered list of distinct primes
\[
 P=(p_1,\ldots,p_n),\qquad n\geq0,
\]
and retain the two archimedean character coordinates $p_+$ and $p_-$.  In the
ordered coordinate basis
\[
 ( [1]_{p_1},\ldots,[1]_{p_n},[p_+],[p_-])
\]
define the already determined direct-sum homomorphism
\begin{equation}
 \Delta_P:\mathbb Z^{n+2}\longrightarrow\mathbb Z^{n+1},
 \qquad
 \Delta_P(k_1,\ldots,k_n,n_+,n_-)
 =(-k_1,\ldots,-k_n,-n_+-n_-).
 \label{eq:global-place-direct-sum-boundary}
\end{equation}
Its first $n$ coordinates are the finite-prime coefficients proved in
Theorem~\ref{thm:primitive-local-boundary}; its last coordinate is the
archimedean map proved in
Theorem~\ref{thm:primitive-archimedean-boundary}.  No place, sign, or
character coordinate has been identified in writing
\eqref{eq:global-place-direct-sum-boundary}.

Let
\begin{equation}
 \Sigma_P:\mathbb Z^{n+1}\longrightarrow\mathbb Z,
 \qquad
 \Sigma_P(y_1,\ldots,y_n,y_\infty)
 =\sum_{i=1}^{n}y_i+y_\infty
 \label{eq:global-place-codiagonal}
\end{equation}
and put
\begin{equation}
 \delta_P^{\mathrm{bal}}=\Sigma_P\circ\Delta_P,
 \qquad
 \delta_P^{\mathrm{bal}}(k_1,\ldots,k_n,n_+,n_-)
 =-\left(\sum_{i=1}^{n}k_i+n_++n_-\right).
 \label{eq:global-place-balancing-row}
\end{equation}
At this point $\delta_P^{\mathrm{bal}}$ is a completely defined homomorphism of
integer lattices.  It is not yet declared to be the boundary map of a global
$C^*$-extension.

\begin{theorem}[Exact balancing lattice]
\label{thm:global-place-balancing-lattice}
Write $m=n+2$ and let $e_1,\ldots,e_m$ be the retained standard basis of
$\mathbb Z^m$.  Then $\delta_P^{\mathrm{bal}}$ is onto and
\[
 \ker\delta_P^{\mathrm{bal}}
 =\left\{x\in\mathbb Z^m:\sum_{j=1}^{m}x_j=0\right\}.
\]
This kernel has the ordered basis
\[
 \alpha_i=e_i-e_{i+1},\qquad 1\leq i\leq m-1,
\]
and the coefficient inverse is explicit:
\begin{equation}
 x=\sum_{i=1}^{m-1}c_i\alpha_i,
 \qquad
 c_i=\sum_{j=1}^{i}x_j.
 \label{eq:global-place-partial-sum-inverse}
\end{equation}
For the standard Euclidean pairing,
\[
 \langle\alpha_i,\alpha_j\rangle
 =\begin{cases}
 2,&i=j,\\
 -1,&|i-j|=1,\\
 0,&|i-j|>1.
 \end{cases}
\]
Thus the balancing kernel, with all coordinates retained, is the root lattice
$A_{m-1}=A_{n+1}$.  More generally, every integer fibre is the affine lattice
\[
 (\delta_P^{\mathrm{bal}})^{-1}(r)
 =-r e_m+\ker\delta_P^{\mathrm{bal}}.
\]
\end{theorem}

\begin{proof}
The formula in \eqref{eq:global-place-balancing-row} shows that the kernel is
the displayed sum-zero lattice.  It is onto because
$\delta_P^{\mathrm{bal}}(-r e_m)=r$ for every $r\in\mathbb Z$.

Each $\alpha_i$ has coordinate sum zero.  Conversely, let
$x=(x_1,\ldots,x_m)$ have sum zero and define the partial sums $c_i$ by
\eqref{eq:global-place-partial-sum-inverse}.  The first coordinate of
$\sum_i c_i\alpha_i$ is $c_1=x_1$; for $2\leq j\leq m-1$ its $j$th coordinate
is
\[
 c_j-c_{j-1}=x_j;
\]
and its last coordinate is
\[
 -c_{m-1}=-\sum_{j=1}^{m-1}x_j=x_m.
\]
This also proves uniqueness of the coefficients, because the first coordinate
recovers $c_1$ and successive coordinate differences recover every $c_i$.
Direct multiplication of the coordinate vectors gives the displayed Cartan
Gram matrix.  Finally, adding $r e_m$ to an element of the fibre over $r$
puts it in the kernel, proving the exact affine-fibre formula.
\end{proof}

For one finite prime, the theorem gives
\[
 \ker\bigl(-(1,1,1):\mathbb Z^3\to\mathbb Z\bigr)
 =\mathbb Z(1,-1,0)\oplus\mathbb Z(0,1,-1)\cong A_2.
\]
This was initially the conjectural stopping point.  The missing codiagonal can,
however, be constructed at the level of stable $C^*$-extensions.  We give the
construction rather than infer it from the displayed integer row.

Put
\[
 \mathcal K=\mathcal K\bigl(\ell^2(\mathbb N_0)\bigr),
 \qquad J=C_0(\mathbb R)\otimes\mathcal K,
\]
and let $e_{00}\in\mathcal K$ be the projection onto the first basis vector.
If
\[
 c(t)=\frac{t-i}{t+i},
\]
then the declared positive generator of $K_1(J)$ is represented in the
unitization by
\begin{equation}
 \gamma(t)=1+(c(t)-1)e_{00}.
 \label{eq:global-place-common-generator}
\end{equation}

\begin{lemma}[A common oriented stabilization]
\label{lem:global-place-common-stabilization}
For every finite prime $p$ and for the archimedean place $\infty$, the proved
local extension admits a stabilized transport
\begin{equation}
 0\longrightarrow J\longrightarrow\widehat E_v
 \longrightarrow\widehat A_v\longrightarrow0,
 \qquad v\in\{p,\infty\},
 \label{eq:global-place-stabilized-local-extension}
\end{equation}
and a morphism from the original local extension to
\eqref{eq:global-place-stabilized-local-extension}.  The quotient algebras are
\[
 \widehat A_p=Q_p\otimes\mathcal K,
 \qquad
 \widehat A_\infty=D_2\otimes\mathcal K.
\]
Under the corner maps on the quotients and the generator
\eqref{eq:global-place-common-generator}, their connecting maps remain
\[
 [1]_p\longmapsto-[\gamma],
 \qquad
 (a[p_+]+b[p_-])\longmapsto-(a+b)[\gamma].
\]
Every stabilization map, generator, and possible orientation reversal is
retained in this assertion.
\end{lemma}

\begin{proof}
For a finite prime, the extension-level Green--Rieffel module in the preceding
subsection gives $I_p\sim_M C_0(\mathbb R)$.  Both algebras are separable and
therefore have countable approximate identities.  The Brown--Green--Rieffel
stabilization theorem gives
\[
 I_p\otimes\mathcal K\cong C_0(\mathbb R)\otimes\mathcal K=J
\]
\cite[Theorems~1.1--1.2, pp.~349--352]{brownGreenRieffel1977}; equivalently,
it follows by taking compact operators after Kasparov stabilization
\cite[Theorem~13.6.2 and Exercise~13.7.1, pp.~131--132]{blackadar1986}.
At the archimedean place, the ideal is
$C_0((0,\infty),M_2)$; the logarithm $(0,\infty)\to\mathbb R$ and a fixed
unitary $\mathbb C^2\otimes\ell^2(\mathbb N_0)\cong
\ell^2(\mathbb N_0)$ give the required stable isomorphism explicitly.

Write $B_v$ for the original local ideal and choose an initial isomorphism
$\theta_v^0:B_v\otimes\mathcal K\to J$.  Its map on the already declared
$K_1$ coordinate is multiplication by a recorded sign
$\varepsilon_v\in\{1,-1\}$.  If $\varepsilon_v=-1$, compose $\theta_v^0$
with the explicit reflection automorphism
\[
 \rho:J\longrightarrow J,\qquad (\rho f)(t)=f(-t),
\]
which takes $[\gamma]$ to $-[\gamma]$; if $\varepsilon_v=1$, leave it
unchanged.  Denote the resulting orientation-preserving isomorphism by
$\theta_v$.  This records the sign correction instead of silently discarding
it.

Tensor the original exact sequence with the nuclear algebra $\mathcal K$.
Let $\tau_v^{\mathrm s}$ be its Busby invariant.  The isomorphism $\theta_v$
extends to multiplier algebras and induces
$\dot\theta_v:Q(B_v\otimes\mathcal K)\to Q(J)$.  Define
\[
 \widehat\tau_v=\dot\theta_v\circ\tau_v^{\mathrm s}
 :\widehat A_v\longrightarrow Q(J)
\]
and let $\widehat E_v$ be its Busby pullback.  The pullback reconstruction and
the induced extension morphism are the constructions of
Sections~15.2--15.3 of Blackadar~\cite[pp.~144--145]{blackadar1986}.  Composing
with the original corner map $x\mapsto x\otimes e_{00}$ gives the asserted
morphism from the unstabilized extension.  Stability sends each displayed
quotient generator to its corresponding rank-one corner, while the declared
$\theta_v$ sends the ideal generator to $[\gamma]$; hence the two local
boundary formulas, including their minus signs, are unchanged.
\end{proof}

Let $\mathcal V=\{\infty\}\cup\{\text{rational primes}\}$ and fix once and for
all a bijection $\nu:\mathcal V\to\mathbb N_0$.  For
\[
 \kappa(j,k)=\frac{(j+k)(j+k+1)}2+k
\]
define an isometry $S_j$ on $\ell^2(\mathbb N_0)$ by
$S_je_k=e_{\kappa(j,k)}$.  Since $\kappa$ is a bijection
$\mathbb N_0^2\to\mathbb N_0$,
\begin{equation}
 S_j^*S_l=\delta_{jl}1,
 \qquad
 \sum_{j=0}^{\infty}S_jS_j^*=1
 \quad\text{strongly}.
 \label{eq:global-place-cuntz-family}
\end{equation}
Write $s_v=1\otimes S_{\nu(v)}\in M(J)$.

\begin{theorem}[Stable place-balancing extension]
\label{thm:global-place-stable-extension}
For every finite prime list $P=(p_1,\ldots,p_n)$ there is an explicitly
constructed extension
\begin{equation}
 0\longrightarrow J\longrightarrow\mathcal E_P
 \longrightarrow\mathcal A_P\longrightarrow0,
 \qquad
 \mathcal A_P=
 \left(\bigoplus_{i=1}^{n}\widehat A_{p_i}\right)
 \oplus\widehat A_\infty,
 \label{eq:global-place-stable-extension}
\end{equation}
with a morphism from the direct sum of the original local extensions after the
declared corner stabilizations.  It has all of the following exact properties.
\begin{enumerate}
\item Its quotient map retains the ordered classes
\[
 ([1]_{p_1},\ldots,[1]_{p_n},[p_+],[p_-])
\]
without merging primes or the two archimedean characters.
\item The ideal map from the direct sum of the stabilized local ideals induces
exactly the codiagonal $\Sigma_P$ of
\eqref{eq:global-place-codiagonal} on $K_1$.
\item The connecting map of \eqref{eq:global-place-stable-extension} is
\[
 \partial_P(k_1,\ldots,k_n,n_+,n_-)
 =-\left(\sum_{i=1}^{n}k_i+n_++n_-\right),
\]
with the displayed sign and all source coordinates retained.  Hence its kernel
is the $A_{n+1}$ lattice of
Theorem~\ref{thm:global-place-balancing-lattice}.
\item Reordering $P$ merely permutes the named quotient summands.  If
$P\subset P'$, zero-inclusion of the new quotient coordinates lifts to a
morphism from the $P$ extension to the $P'$ extension which is the identity on
$J$.  Thus adjoining a prime and removing it by restriction to the zero
coordinate are compatible at extension level, not only on $K$-groups.
\end{enumerate}
\end{theorem}

\begin{proof}
Put $V_P=\{p_1,\ldots,p_n,\infty\}$ and form the finite direct sums
\[
 B_P=\bigoplus_{v\in V_P}J,
 \qquad
 \widehat E_P^\oplus=\bigoplus_{v\in V_P}\widehat E_v,
 \qquad
 \mathcal A_P=\bigoplus_{v\in V_P}\widehat A_v.
\]
For $(b_v)_{v\in V_P}\in B_P$, define
\begin{equation}
 \mu_P((b_v)_v)=\sum_{v\in V_P}s_vb_vs_v^*\in J.
 \label{eq:global-place-ideal-map}
\end{equation}
The orthogonality in \eqref{eq:global-place-cuntz-family} gives, coordinate by
coordinate,
\[
 \mu_P(b)\mu_P(d)
 =\sum_{v,w}s_vb_vs_v^*s_wd_ws_w^*
 =\sum_vs_vb_vd_vs_v^*=\mu_P(bd),
\]
and $s_v^*\mu_P(b)s_v=b_v$.  Thus $\mu_P$ is an injective
$*$-homomorphism.  The same finite formula defines
\[
 \overline\mu_P:M(B_P)=\bigoplus_{v\in V_P}M(J)\longrightarrow M(J).
\]
It maps $B_P$ into $J$ and therefore induces a corona homomorphism
$\dot\mu_P:Q(B_P)\to Q(J)$.

Let
\[
 \widehat\tau_P^\oplus=\bigoplus_{v\in V_P}\widehat\tau_v:
 \mathcal A_P\longrightarrow Q(B_P)
\]
be the Busby invariant of the direct-sum extension and set
\begin{equation}
 \tau_P=\dot\mu_P\circ\widehat\tau_P^\oplus.
 \label{eq:global-place-busby-sum}
\end{equation}
The required middle algebra is the explicit pullback
\begin{equation}
 \mathcal E_P=
 \{(a,m)\in\mathcal A_P\oplus M(J):
       \tau_P(a)=\pi_J(m)\},
 \label{eq:global-place-busby-pullback}
\end{equation}
where $\pi_J:M(J)\to Q(J)$ is the quotient map.  Projection onto $a$ is onto,
and its kernel is exactly $\{(0,j):j\in J\}$, proving
\eqref{eq:global-place-stable-extension}.  In the Busby pullback models for the
local extensions, the middle map is
\begin{equation}
 ((a_v,m_v))_{v\in V_P}
 \longmapsto
 \left((a_v)_{v\in V_P},
       \sum_{v\in V_P}s_vm_vs_v^*\right).
 \label{eq:global-place-middle-map}
\end{equation}
Together with $\mu_P$ on ideals and the identity on $\mathcal A_P$, this is a
morphism of short exact sequences.  It is an extension-level construction
before any $K$-group calculation.  It is the finite orthogonal-corner form of
the standard Busby sum for stable ideals
\cite[Sections~15.6 and 15.9, pp.~149--150, 156]{blackadar1986}.

It remains to compute the ideal map without appealing to its intended answer.
The $v$th copy of \eqref{eq:global-place-common-generator} is sent to
\[
 1+(c(t)-1)q_v,
 \qquad q_v=S_{\nu(v)}e_{00}S_{\nu(v)}^*.
\]
The projection $q_v$ is rank one.  A rotation in the at most two-dimensional
span of $e_0$ and $S_{\nu(v)}e_0$ gives an explicit norm-continuous path of
rank-one projections from $e_{00}$ to $q_v$.  Substituting that path for
$e_{00}$ in \eqref{eq:global-place-common-generator} proves that the $v$th
corner map takes $[\gamma]$ to $[\gamma]$.  Since
$K_1(B_P)=\bigoplus_{v\in V_P}\mathbb Z[\gamma]$, additivity gives
\[
 (\mu_P)_*((y_v)_v)=\left(\sum_{v\in V_P}y_v\right)[\gamma].
\]
This is exactly $\Sigma_P$, not an unidentified isomorphism.

The top extension in the morphism diagram has boundary
$\Delta_P$ from \eqref{eq:global-place-direct-sum-boundary}.  The connecting
maps in $K$-theory are natural for morphisms of short exact sequences
\cite[Definition~21.1.1 and Example~21.1.2(a), pp.~266, 268]{blackadar1986}.
Consequently
\[
 \partial_P=(\mu_P)_*\circ\Delta_P
 =\Sigma_P\circ\Delta_P=\delta_P^{\mathrm{bal}},
\]
which proves the first three assertions and invokes the already proved exact
lattice theorem for the kernel.

Finally, the operators $s_v$ and the transports $\theta_v$ were fixed by the
place $v$, not by its position in $P$.  Reordering therefore only permutes the
direct-sum notation.  If $P\subset P'$, let
$\iota_{P,P'}:\mathcal A_P\to\mathcal A_{P'}$ insert zero in the new place
coordinates.  The literal formulas give
\[
 \tau_{P'}\circ\iota_{P,P'}=\tau_P.
\]
Hence $(a,m)\mapsto(\iota_{P,P'}(a),m)$ maps
$\mathcal E_P$ to $\mathcal E_{P'}$, is the identity on $J$, and makes the
extension diagram commute.  Restriction back to the zero-coordinate summand
recovers the $P$ extension.  This proves the fourth assertion.
\end{proof}

\paragraph{Passage to the natural arithmetic representative.}
The extension-existence, codiagonal, sign, kernel, permutation, and finite-set
compatibility assertions are therefore theorems after stabilization.  The
audited source itself does not define its larger symbol $A_{\varnothing p}$ or
a simultaneous multi-place action.  Subsection
\ref{subsec:unstabilized-semilocal-one-zero} therefore constructs the natural
unit-retaining semilocal crossed product independently and proves its exact
$K$-theory comparison with this stable scalar extension.  Subsection
\ref{subsec:stable-to-natural-extension} then constructs the formerly missing
extension-level $*$-homomorphism by an exact multiplier/corona Busby equation.
The natural boundary still collapses the two archimedean character classes in
$K_0$, and the natural ideal still contains nonconstant unit-fibre classes;
the new theorem records these as the kernel and image data of an injective
stable extension morphism.  No endpoint statement about zeta zeros follows
from any of these extension theorems.

\paragraph{Literature and novelty boundary.}
The proof uses the content-read stable-isomorphism theorem of Brown, Green, and
Rieffel.  It also uses Blackadar's Busby pullback, stable extension sum and
functoriality, and the naturality of $K$-theory connecting maps.  A bounded canonical-index query
for these exact dependencies routed the cited sources.  The earlier bounded
query for the combined terms ``codiagonal'', ``$A_2$'', ``boundary'', and
``$K$-theory'', followed by a phrase-level arXiv check, did not locate this
place-indexed diagram.  That search scope supports no general novelty claim.
The theorem above is the explicit stable construction.  The next two
subsections supply the independently defined arithmetic target and the exact
stable-to-natural extension morphism.  The bounded search scope supports no
general novelty claim.

\paragraph{Formal finite-coordinate replay.}
The file
\begin{center}
\footnotesize
\path{formal/lean/ZetaFunctionFoundation/GlobalPlaceBalancing.lean}
\end{center}
certifies the direct-sum boundary, codiagonal composite, scalar-kernel equation,
surjectivity, and the naturality-implied formula for an arbitrary finite number
of prime coordinates.  Receipt
\path{formal/lean/receipts/FORMAL-20260830-0036.receipt.json} records Lean
4.32.0, Mathlib 4.32.0, exit code zero, and a peak aggregate working set below
the authorized five-gigabyte ceiling.  The Busby invariant, multiplier-algebra
construction, and $C^*$-algebraic naturality theorem remain certified by the
complete standard proof above and its human literature dependencies; this
Lean file does not pretend to formalize those analytic objects.
